\documentclass[12pt]{article} \usepackage{mathtools}
\usepackage{amssymb}
\usepackage{fontspec}
\usepackage{algorithm}
\usepackage{algpseudocode}
% \usepackage[letf=1.06cm,top=0.9cm,right=1.06cm,bottom=0.49cm]{geometry}
% \setmainfont{GFS Didot}
\setmainfont{EB Garamond}


% Set page size and margins
\usepackage[a4paper,top=1.5cm,bottom=1.5cm,left=1.3cm,right=1.3cm]{geometry}
\usepackage{fancyhdr}
\pagestyle{fancy}
\setlength{\parindent}{0pt}

% Useful packages
\usepackage{pdflscape}
\usepackage{minted}
\usepackage{graphicx}
\graphicspath{{./media/}}
\usepackage{subfig}
\usepackage{float}
\usepackage[colorlinks=true, allcolors=blue]{hyperref}
\usepackage{siunitx}
\usepackage{bm}
% \usepackage{sectsty}
% \sectionfont{\fontsize{12}{15}\selectfont}

\title{\vspace{-2cm}096 - Ειδικές Κεραίες\\
    \large Εργασία 2η \\
    Αλγόριθμοι Μορφοποίησης Δέσμης Στοιχειοκεραιών \\
    NSB}

\author{Ιωάννης Δημουλιός \\ 10641}
\date{Εαρινό εξάμηνο 2026}

\lhead{096 - Ειδικές Κεραίες}
\chead{Εργασία 2η}
\rhead{Ιωάννης Δημουλιός}
\begin{document}
\maketitle
\textbf{Σχόλια για την υλοποίηση}\\
Ο πίνακας $A^H A$ σε αρκετούς από τους απαιτούμενος υπολογισμούς ήταν ιδιάζων, οπότε για να τον αξιοποιήσουμε στον υπολογισμό των βαρών του NSB χρησιμοποιούμε τον τύπο
\[
w = A(A^H A + \sigma^2 I)^{-1} e_1
,\]
όπου $\sigma^2$ η διακύμανση του θορύβου. \\
Για την επιλογή των γωνιών άφιξης πλασματικών σημάτων παρεμβολής χρησιμοποιήθηκε ο εξής απλός άπληστος αλγόριθμος. \\
Κατασκευάζουμε το διάγραμμα ακτινοβολίας και τοποθετούμε ένα σήμα παρεμβολής στη γωνία $\theta_i$ που εντοπίζεται ο μεγαλύτερος πλευρικός λοβός και υπολογίζουμε εκ νέου τα βάρη NSB. Επαναλαμβάνουμε έως ότου εξαντληθούν τα $18$ διαθέσιμα σήματα παρεμβολής ή ο μεγαλύτερος πλευρικός λοβός είναι κατά τουλάχιστον \SI{20}{\dB} μικρότερος του κύριου λοβού. \\

\textbf{Σχόλια για τα αποτελέσματα} \\
i) Απόκλιση κύριου λοβού \\
Η βελτιστοποίηση φαίνεται να λειτουργεί ικανοποιητικά όταν $\delta \geq 10^\circ$και για όλα τα $\delta$ όταν $\text{SNR}\leq \SI{0}{\dB}$. Πράγματι σε αυτήν την περίπτωση πετυχαίνουμε μικρές αποκλίσεις γωνιών για τον κεντρικό λοβό $|\Delta\theta_0|\leq 0.5^\circ$. Όσο το $\delta$ αυξάνεται, η απόκλιση μειώνεται. \\
Ωστόσο, για $\text{SNR} \geq \SI{10}{\dB}$ και $\delta \leq 8^\circ$, η απόκλιση του κύριου λοβού αποκλίνει αριθμητικά. Αυτό πιστεύουμε ότι οφείλεται στον σχεδόν ιδιάζων πίνακα $A^H A$ και στον μικρό θόρυβο που δεν μπορεί να βελτιώσει το condition number (100K) επαρκώς όταν χρησιμοποιούμε τον τύπο παραπάνω για τον υπολογισμό των βαρών. Ταυτόχρονα για μικρά $\delta$ οι μηδενισμοί είναι πολύ κοντά μεταξύ τους και με την προσθήκη των πλασματικών μηδενισμών σε ορισμένες περιπτώσεις επηρεάζεται και ο κύριος λοβός. Ενδεχομένως, μια διαφορετική τεχνική δίχως εξάρτηση από τον θόρυβο για την αποφυγή ιδάζοντος πίνακα να έδινε καλύτερα αποτελέσματα. \\

ii) Απόκλιση μηδενισμών \\
Οι μηδενισμοί παρουσιάζουν μικρότερα σφάλματα ακόμη και όταν ο απόκλιση του κύριου λοβού εκτοξεύεται. Πράγματι σε όλες τις περιπτώσεις $|\Delta\theta_i| \leq 4.7^\circ$ με τα μεγαλύτερα σφάλματα να παρουσιάζονται για μικρά $\delta$ και SNR. Τα νούμερα ουσιαστικά υποστηρίζουν τον ισχυρισμό ότι οι μηδενισμοί δεν επηρεάζουν άλλους μηδενισμούς, αλλά μόνο τον κύριο λοβό. \\

iii) SINR \\
Το SINR εύλογα είναι αύξουσα συνάρτηση του SNR. Στις περιπτώσεις που η απόκλιση του κύριου λοβού είναι προβληματική η ελάχιστη τιμή του SINR είναι πολύ μικρότερη από τη μέγιστη και τον μέσο όρο και σαφώς η τυπική του απόκλιση είναι αυξημένη. \\
Η μόνη επιπλεόν παρατήρηση είναι ότι για μικρά SNR όταν ο θόρυβος ακόμη κυριαρχεί στον παρανομαστή του SINR, αυτό δεν βελτιώνεται για μεγαλύτερα $\delta$ (βλ. για $\text{SNR} = \SI{-10}{\dB}$ και $\delta \geq 10^\circ$). \\

iv) SLL \\
Η σχεδιαστική προδιαγραφή του $\text{SLL} \leq \SI{-20}{\dB}$ κατά κανόνα επιτυγχάνεται εκτός των περιπτώσεων προβληματικής απόκλισης του κεντρικού λοβού. Για μικρά $\delta$ γενικότερα τo μέγιστο SLL μπορεί να είναι κοντά αλλά να μην πετυχαίνει την τιμή που ορίζει η σχεδιαστική προδιαγραφή.


\clearpage
\newgeometry{
    left=4cm,
    right=3.5cm,
    top=0.5cm,
    bottom=0.5cm
}
\begin{landscape}
\thispagestyle{empty}
\textbf{NSB}\\
\begin{tabular}{rrrrrrrrrrrrrrrrrrr}
\toprule
snr & \delta & \min \Delta\theta_0 & \max \Delta\theta_0 & \overline{\Delta\theta_0} & \text{std} \Delta\theta_0 & \min \Delta\theta_i & \max \Delta\theta_i & \overline{\Delta\theta_i} & \text{std} \Delta\theta_i  & \min\text{sinr} & \max\text{sinr} & \overline{\text{sinr}} & \text{std sinr} & \min\text{sll} & \max\text{sll} & \overline{\text{sll}} & \text{std sll} \\
\midrule
-10 & 6 & -2.100 & 2.100 & -0.000 & 0.341 & -4.700 & 4.700 & -0.014 & 1.212 & 1.119 & 3.493 & 3.271 & 0.393 & -21.851 & -18.845 & -20.554 & 0.452 \\
-10 & 8 & -0.900 & 0.900 & -0.000 & 0.144 & -3.200 & 3.200 & -0.045 & 1.258 & 3.186 & 3.507 & 3.413 & 0.041 & -22.168 & -20.005 & -20.741 & 0.528 \\
-10 & 10 & -0.200 & 0.200 & -0.000 & 0.088 & -2.700 & 2.700 & -0.002 & 0.826 & 3.383 & 3.570 & 3.503 & 0.051 & -21.801 & -20.001 & -20.707 & 0.436 \\
-10 & 12 & -0.200 & 0.200 & -0.000 & 0.061 & -3.300 & 3.100 & -0.034 & 1.153 & 3.412 & 3.535 & 3.466 & 0.027 & -22.112 & -20.000 & -20.692 & 0.480 \\
-10 & 14 & -0.200 & 0.200 & -0.000 & 0.080 & -2.900 & 2.900 & -0.009 & 0.936 & 3.432 & 3.527 & 3.489 & 0.023 & -21.738 & -20.009 & -20.910 & 0.514 \\
-10 & 16 & -0.100 & 0.100 & -0.000 & 0.060 & -3.000 & 3.000 & -0.019 & 1.071 & 3.451 & 3.532 & 3.480 & 0.019 & -21.743 & -20.085 & -20.988 & 0.447 \\
0 & 6 & -4.300 & 4.300 & -0.000 & 0.799 & -4.700 & 4.700 & -0.005 & 1.106 & 5.445 & 13.463 & 11.949 & 1.489 & -22.019 & -14.734 & -20.187 & 1.193 \\
0 & 8 & -2.900 & 2.900 & -0.000 & 0.460 & -3.000 & 3.000 & -0.007 & 0.763 & 10.894 & 13.423 & 13.000 & 0.304 & -26.109 & -19.323 & -21.186 & 1.032 \\
0 & 10 & -1.700 & 1.700 & -0.000 & 0.303 & -2.400 & 2.400 & -0.005 & 0.503 & 12.910 & 13.326 & 13.190 & 0.104 & -24.716 & -20.006 & -21.283 & 0.894 \\
0 & 12 & -0.400 & 0.400 & -0.000 & 0.137 & -3.800 & 3.800 & -0.009 & 0.593 & 13.038 & 13.296 & 13.128 & 0.060 & -23.147 & -20.004 & -21.100 & 0.855 \\
 0 & 14 & -0.500 & 0.500 & -0.000 & 0.160 & -2.500 & 2.200 & -0.009 & 0.498 & 13.057 & 13.289 & 13.165 & 0.048 & -25.588 & -20.003 & -21.657 & 1.738 \\
 0 & 16 & -0.200 & 0.200 & -0.000 & 0.117 & -2.200 & 2.200 & -0.004 & 0.447 & 13.108 & 13.261 & 13.166 & 0.028 & -24.605 & -20.065 & -21.803 & 1.237 \\
 10 & 6 & -124.300 & 124.300 & -0.000 & 12.813 & -3.900 & 3.900 & -0.005 & 0.757 & 9.560 & 23.732 & 20.898 & 3.234 & -21.220 & -5.170 & -16.640 & 3.351 \\
 10 & 8 & -5.000 & 5.000 & -0.000 & 0.734 & -3.600 & 3.600 & -0.004 & 0.431 & 16.345 & 23.656 & 22.710 & 1.072 & -26.484 & -11.514 & -20.680 & 2.112 \\
 10 & 10 & -2.300 & 2.300 & -0.000 & 0.276 & -1.100 & 1.100 & -0.001 & 0.180 & 22.275 & 23.335 & 23.091 & 0.165 & -23.524 & -17.401 & -21.164 & 0.962 \\
 10 & 12 & -0.400 & 0.400 & -0.000 & 0.148 & -1.500 & 1.500 & -0.002 & 0.160 & 22.828 & 23.240 & 23.037 & 0.079 & -23.572 & -20.023 & -21.302 & 0.942 \\
 10 & 14 & -0.400 & 0.400 & -0.000 & 0.161 & -1.200 & 1.200 & -0.001 & 0.131 & 22.908 & 23.241 & 23.084 & 0.069 & -26.031 & -20.001 & -21.719 & 1.745 \\
 10 & 16 & -0.300 & 0.300 & -0.000 & 0.117 & -0.900 & 0.900 & -0.002 & 0.127 & 23.019 & 23.209 & 23.101 & 0.034 & -24.747 & -20.047 & -21.943 & 1.220 \\
 20 & 6 & -133.300 & 133.300 & -0.000 & 20.430 & -3.100 & 3.100 & -0.001 & 0.224 & 7.922 & 33.733 & 31.640 & 5.707 & -20.634 & -7.040 & -15.877 & 2.897 \\
 20 & 8 & -110.500 & 110.500 & -0.000 & 9.813 & -1.800 & 1.800 & -0.001 & 0.099 & 8.590 & 33.720 & 32.809 & 2.140 & -26.367 & -11.507 & -20.480 & 2.502 \\
 20 & 10 & -2.400 & 2.400 & -0.000 & 0.242 & -0.500 & 0.400 & -0.000 & 0.039 & 32.610 & 33.332 & 33.079 & 0.169 & -23.552 & -17.109 & -21.161 & 1.009 \\
 20 & 12 & -0.400 & 0.400 & -0.000 & 0.143 & -0.700 & 0.700 & -0.000 & 0.034 & 32.788 & 33.235 & 33.026 & 0.083 & -23.623 & -20.015 & -21.289 & 0.962 \\
 20 & 14 & -0.400 & 0.400 & -0.000 & 0.158 & -0.200 & 0.200 & -0.001 & 0.020 & 32.832 & 33.235 & 33.071 & 0.074 & -26.061 & -20.028 & -21.847 & 1.722 \\
 20 & 16 & -0.300 & 0.300 & -0.000 & 0.116 & -0.300 & 0.300 & -0.001 & 0.025 & 33.008 & 33.203 & 33.093 & 0.035 & -24.765 & -20.011 & -21.944 & 1.216 \\
\bottomrule
\end{tabular}
\end{landscape}
\restoregeometry
\end{document}
